\chapter{Calculated columns}\label{ch:calc}
\begin{goals}
\textbf{In this chapter} you derive new columns from a formula, check them before adding them, manage
them, and learn what they cost and where they go.
\end{goals}

A \term{calculated column} is a column lamina works out for each row from the row's other columns:
\code{dx * dy * dz * density} for a block's tonnes, \code{if(Cu\_pct >= 0.3, "ore", "waste")} for a class,
\code{bin(Z, 15)} for a bench. It is written in the expression language (\cref{ch:expr}), it is worked out
as each row is read, and it is never written to your file. Once added it is a column like any other:
the grid shows it, and filters, sorts, statistics, group by, grade–tonnage and export all use it.

\section{Adding one}\label{sec:add-calc}

Click \ui{ƒ+ col} on the toolbar, or choose \ui{Add calculated column…} on any column header's menu.
The editor asks for two things:

\begin{description}
\item[\ui{Name}] a letter, then letters, digits or \code{\_}. It can't be the name of a column already
in the table. (No hyphen: in an expression a hyphen is a minus, \cref{sec:names}.)
\item[\ui{Expr}] the expression. It is coloured and completed as the filter box is
(\cref{sec:filter-box}).
\end{description}

As you type, the editor checks the expression and says what it uses, \code{✓ uses: dx, dy, dz, density},
or what is wrong with it, and previews its value on the first six rows of the file, so you can check the
arithmetic on rows you can see. \ui{Add}, or \key{Enter} in the expression, adds the column; \key{Esc} or
\ui{Cancel} closes the editor.

\shot[0.6\linewidth]{calc-editor}{The calculated-column editor: a name, the expression, what it uses, and
its value on the first six rows.}

The new column appears at the right of the table, its header marked \ui{ƒ}.

\subsection{The column's type}\label{sec:calc-type}

lamina decides whether the new column holds numbers or text by working it out on the sample spread
through the file (\cref{sec:gutter-sample}): a number column when every non-blank result there is a
number (or text that reads as one), a text column otherwise. Spread, so that an expression blank on the
file's first rows (\code{if(Cu\_pct >= 1, Cu\_pct, blank)} on a model whose first rows are waste) is still
seen to give numbers further down. If it decides wrong, the header's menu has \ui{Treat as number} and
\ui{Treat as text}.

\section{Using another calculated column}

An expression may use a calculated column added before it: add \col{tonnes} as \code{dx * dy * dz *
density}, then \col{cu\_t} as \code{tonnes * Cu\_pct / 100}. Each row's values are worked out in the order
the columns were added.

\section{Editing and removing}\label{sec:calc-manager}

A calculated column's header menu has \ui{Edit …} and \ui{Remove …}. \menu{Data>Calculated columns…}
lists them all, with their expressions, to edit, remove or add from one place. Editing a column the
filter or another column uses changes what they show; removing one they use leaves the filter or the
column referring to a name that isn't there, which lamina reports as for any unknown column.

\section{What they cost}\label{sec:calc-cost}

A calculated column is worked out every time a row is read: on screen, and in every scan. The arithmetic
itself is fast; what can cost is that it may need columns a scan otherwise wouldn't read. On a
\file{.dm} or a Parquet file, a filter normally reads only the columns it names (\cref{sec:strided}). With
any calculated column defined, it reads every column, because lamina doesn't follow which columns a
filter reaches through a calculated one. On delimited text, which is read whole anyway, it makes no
difference.

\section{Where they go}\label{sec:calc-where}

\begin{itemize}
\item \textbf{Export} writes calculated columns as ordinary columns of values (\cref{ch:export}). A
\file{.dm} with a calculated tonnes column exported to CSV gives a CSV with a real tonnes column.
\item \textbf{A lens} keeps their names and expressions, and adds them to the next file it is applied to
(\cref{ch:lens}).
\item \textbf{Closing the view}, or lamina, drops them, as it drops the filter and the sort. They were
never in the file.
\end{itemize}

\section{Bands to group by}\label{sec:bands}

\code{bin} turns a number into the lower edge of its band, so a calculated column of it gives categories
to group by: \code{bin(Z, 15)} for 15~m benches, \code{bin(Cu\_pct, 0.25)} for grade classes,
\code{bin(Y, 100)} for 100~m slices north. Group by the band (\cref{ch:groupby}) and you get a table by
bench, by class, by slice: a swath plot's numbers, without leaving the table.

\begin{tryit}
Tonnes and copper metal per bench.
\begin{enumerate}
\item In the sample, add \col{t} with \code{dx * dy * dz * density * ore\_pct}: the ore tonnes in each
block. The preview shows 0 for the first rows, which are waste.
\item Add \col{cu\_t} with \code{t * Cu\_pct / 100}: tonnes of copper.
\item Add \col{bench} with \code{bin(Z, 20)}.
\item Open \menu{Tools>Group by…}, group by \col{bench}, with values \col{t} and \col{cu\_t}. The table
gives, per 20~m bench, the sum of each: ore tonnes and contained copper.
\end{enumerate}
\end{tryit}
